\section{整数分拆：重数限制的差一问题}\index{integer partitions}\index{pentagonal number theorem}
实现见第~\pageref{compact-Partitions}~页。分拆不区分加数顺序，部分均为正整数。
无约束生成函数 $P(z)=\prod_{i\ge1}(1-z^i)^{-1}$，空分拆给出 $p(0)=1$。
五边形数定理为
\[
\prod_{i\ge1}(1-z^i)=1+\sum_{j\ge1}(-1)^j
\left(z^{j(3j-1)/2}+z^{j(3j+1)/2}\right).
\]
与 $P(z)$ 相乘后比较系数，得到 $p(n)$ 的递推，符号按 $+,+,-,-,+,+,\ldots$ 成对出现。
每项只依赖更小的下标，每个 $n$ 有 $O(\sqrt n)$ 个非零偏移。
实现先用 64 位累加，再归一化到模数范围；无需逆元，合数模数与模数 1 均可使用。

每个正整数最多出现 $k$ 次的生成函数是
\[
\prod_{i\ge1}(1+z^i+\cdots+z^{ki})
=P(z)\prod_{i\ge1}(1-z^{(k+1)i}).
\]
因此用已有 $p$ 表与放大 $k+1$ 倍的五边形偏移相乘，符号为 $-,-,+,+,\ldots$。
本库 \texttt{limited(n,k)} 的上限包含 $k$；$k=0$ 时非零 $n$ 无解，$k=1$ 为互异分拆。
kuangbin 的注释写“不超过 k 个”，但其 \texttt{solve(n,k)} 使用放大 $k$ 倍的公式，
实际表示每种数少于 $k$ 次，对应本库 \texttt{limited(n,k-1)}（$k\ge1$）。
HDU 原题页面本次无法读取，不能由这条注释推断原题措辞。

两套实现以任意精度完全背包验证到 1000，以逐种部分、逐个重数枚举验证有界情况到 100，
覆盖素数与合数模数、0/1 重数限制和大于总和的限制；另检查 $N=100000$ 的边界。
普通与 ASan/UBSan 均通过，暂无在线 AC。


\Needspace{9\baselineskip}
Library Checker 的分拆函数要求输出 $p(0),\ldots,p(N)$，$N\le500000$，
完整用法见第~\pageref{usage-example-158}~页；固定官方11组数据的本地checker已通过。
P6189 将非增跑步距离序列对应到无序分拆，模数不保证素数，
用法见第~\pageref{usage-example-159}~页。它是比赛应用，不替代正式模板题。
额外独立参考将部分分成不超过 $B$ 与大于 $B$ 两组：前者完全背包，
后者按恰好 $j$ 个部分维护 $g_j(s)=g_j(s-j)+g_{j-1}(s-B-1)$，再卷积。
取 $B=\lfloor\sqrt N\rfloor$，时间 $O(N\sqrt N)$、滚动数组空间 $O(N)$，
不复用被测的五边形数递推；覆盖 $N=10^5$ 的合数模数与模数1。
这两份用法只调用构造和 $p$ 表，不代表 limited 接口已有正式题的线上验证。

来源：kuangbin 2.11；\url{https://oi-wiki.org/math/combinatorics/partition/}。

